\documentclass[11pt,letterpaper]{article}
\usepackage[utf8]{inputenc}
\usepackage{csquotes}
\usepackage[spanish]{babel}
\usepackage{geometry}
\usepackage{fancyhdr}
\usepackage{booktabs}
\usepackage{longtable}
\usepackage{array}
\usepackage{enumitem}
\usepackage{tikz}
\usepackage{graphicx}
\usepackage{hyperref}
\usepackage{ragged2e}
\usepackage{amsmath}
\usepackage{pdflscape}
\usepackage{tabularx}
\usepackage{caption}
\usepackage{siunitx}
\usepackage[spanish]{cleveref}
\usepackage{microtype}
\usepackage[
backend=biber,
style=authoryear,
sorting=nyt
]{biblatex}

\addbibresource{references.bib}

\geometry{
  a4paper,
  left=1in,
  right=1in,
  top=1in,
  bottom=1in
}

\pagestyle{fancy}
\fancyhf{}
\fancyhead[L]{\textit{Diagnostico}}
\fancyfoot[C]{\thepage}


\begin{document}
\begin{center}

  \vspace*{2cm}

  \textbf{Análisis Operativo y Estratégico del Relleno Sanitario Doña Juana}

  \vspace{0.5cm}

  \textbf{Diagnóstico}

  \vspace{3cm}

  Universidad Distrital Francisco José de Caldas

  \vspace{0.5cm}

  Programa Académico de Física

  \vspace{2cm}

  Gestión Tecnologica y Compromiso Social

  \vspace{0.5cm}

  Adolfo Leon Agaton

  \vspace{3cm}

  Julian Avila \\
  Bryan Martinez \\
  Angie Sierra

  \vspace{1.5cm}

  \today

\end{center}

\newpage

\tableofcontents
\newpage

\section{Contexto Estratégico Ejecutivo y Anatomía Macroambiental}

La gestión, mitigación y valorización de los residuos sólidos urbanos constituye
un problema infraestructural y ambiental fundamental para las megaciudades. En
la sabana de Bogotá, Colombia, el Relleno Sanitario Doña Juana ha operado como
el repositorio exclusivo de desechos metropolitanos desde 1988. Situado en la
localidad de Ciudad Bolívar, abarca 623 hectáreas entre Mochuelo Alto y Mochuelo
Bajo. Recibe más de 2.2 millones de toneladas anuales de residuos de forma
continua, acumulando una fracción orgánica que representa el 52\% de la masa
total.

La descomposición anaeróbica de esta fracción orgánica genera biogás, una mezcla
conformada principalmente por metano ($\text{CH}_4$) y dióxido de carbono
($\text{CO}_2$). El metano es un gas de efecto invernadero con un potencial de
calentamiento global 28 veces superior al del dióxido de carbono en un horizonte
de cien años. Sin estrategias de mitigación, la emisión de este gas contribuye
al cambio climático antropogénico y provoca problemas de salud pública local,
contaminación por olores, proliferación de vectores e inestabilidad geotécnica.

En este entorno, Biogás Colombia S.A.S. E.S.P., operando como Biogás Doña Juana
S.A.S. E.S.P. y representada por Helmuth Gallego, actúa como la concesionaria
responsable de la mitigación ambiental y valorización energética del biogás.
Mediante una concesión de 35 años adjudicada en 2007 por el Distrito Capital y
la Unidad Administrativa Especial de Servicios Públicos (UAESP), la entidad
procesa y monetiza el biogás, consolidándose como un proyecto de economía
circular en Colombia.

La infraestructura física requiere ingeniería dinámica. El biogás se extrae
mediante 290 pozos verticales distribuidos en las celdas de residuos. Estos
pozos se conectan a través de 76 kilómetros de tuberías de polietileno de alta
densidad, las cuales conducen el gas a una planta central. Sistemas de
adquisición de datos y control de supervisión regulan la presión y la dinámica
de fluidos para mantener una tasa de extracción entre 11,700 y 16,000 metros
cúbicos por hora.

En la planta central, el gas es impulsado por cuatro sopladores industriales
hacia tres antorchas cerradas. Cada chimenea tiene una capacidad de 5,000 metros
cúbicos normales por hora ($\text{Nm}^3/\text{h}$), donde el metano entra en
combustión a temperaturas superiores a 500 grados Celsius. Esta destrucción
térmica evita la emisión de 800,000 toneladas de dióxido de carbono equivalente
($\text{tCO}_2\text{e}$) anuales.

El objetivo estratégico de la empresa es la transición de la quema térmica a la
generación eléctrica a gran escala y la inyección de biometano. Esta transición
representa una actualización de capacidad actualmente limitada por cuellos de
botella estructurales y financieros.

\newpage

\section{Identidad Corporativa y Arquitectura Organizacional}

La estructura jurídica y operativa de la entidad responde a la complejidad de
los mercados regulados de servicios públicos y los mecanismos internacionales de
financiación climática.

\subsection{Naturaleza Jurídica y Tipo de Entidad}

La organización opera bajo la razón social \textbf{Biogás Doña Juana S.A.S.
E.S.P.}, siendo reconocida comercialmente en el sector como Biogás Colombia.
Jurídicamente, se constituye como una Empresa de Servicios Públicos (E.S.P.) de
carácter privado. Esta naturaleza le otorga un mandato especializado y distinto
al del operador del relleno (Centro de Gerenciamiento de Residuos - CGR),
enfocándose exclusivamente en la captura y valorización del componente gaseoso.
Aunque es una corporación de derecho privado, su modelo de ingresos y operación
mantiene un vínculo indisoluble con el interés público a través de su contrato
de concesión y el esquema de regalías transferidas a la UAESP.

\subsection{Evolución de la Estructura de Gobernanza}

La arquitectura corporativa actual es el resultado de una evolución forzada por
la volatilidad de los precios del carbono. Originalmente estructurada para
apalancar los ingresos del Mecanismo de Desarrollo Limpio (MDL) con
participación de actores internacionales como Gas Natural SDG S.A. y Proactiva
Medio Ambiente, la empresa debió someterse a una reorganización obligatoria bajo
la Ley 1116 tras el colapso de los precios de las Reducciones Certificadas de
Emisiones (CER) en 2011.

La gobernanza actual emerge de la exitosa negociación de este proceso en 2014,
lo que permitió a la entidad desacoplar su supervivencia operativa de la
volatilidad del mercado spot de carbono y establecerse como una Sociedad por
Acciones Simplificada (S.A.S.) estabilizada. Esta estructura reorganizada ha
dotado a la empresa de la capacidad jurídica y financiera para ejecutar
proyectos de capital a largo plazo, como la actual expansión de biometano.

\subsection{Estructura Organizacional y Fuerza Laboral Estimada}

Dada la envergadura de la concesión, que abarca 623 hectáreas y opera 290 pozos
de extracción de manera continua, la arquitectura organizacional requiere un
equipo técnico y administrativo altamente especializado. Se estima que la planta
de personal directo de Biogás Doña Juana S.A.S. E.S.P. se compone de
aproximadamente 45 a 50 trabajadores fijos, estructurados para garantizar una
operatividad de 8760 horas al año (24/7).

Para sostener la generación continua y el control del campo de biogás, el área
operativa se divide en un sistema de cuatro turnos rotativos. La distribución de
la fuerza laboral se desglosa de la siguiente manera:

\begin{itemize}

  \item \textbf{Operación de Planta y GCCS (15-20 personas):} Incluye ingenieros
    de turno, operadores de consola (SCADA) y técnicos de campo encargados de la
    calibración de cabezales y purga de condensados en los 76 km de tubería.

  \item \textbf{Mantenimiento Electromecánico (10-12 personas):} Personal
    especializado en mecánica térmica y electricidad de potencia, enfocados en
    el mantenimiento preventivo (cambio de aceite, bujías, ajuste de válvulas)
    de la flota de motogeneradores Jenbacher y los sopladores industriales.

  \item \textbf{Ingeniería Química y Acondicionamiento (5-7 personas):}
    Analistas de laboratorio y técnicos encargados del monitoreo de las unidades
    de remoción de siloxanos, control de humedad y cromatografía del gas.

  \item \textbf{Administración, HSE y Sostenibilidad (10-11 personas):}
    Gerencia, especialistas en mercados de carbono (comercialización de CER),
    contabilidad, y gestores de relaciones comunitarias y monitoreo de emisiones
    ambientales.

\end{itemize}

\begin{figure}[htbp] \centering \resizebox{\textwidth}{!}{
    \begin{tikzpicture}[ level 1/.style={sibling distance=8.5cm}, level
      2/.style={sibling distance=4.2cm}, level distance=2.5cm, every
      node/.style = { shape=rectangle, rounded corners, draw, align=center,
        top color=white, bottom color=gray!10, font=\footnotesize, minimum
      height=1.2cm, text width=3.8cm }, edge from parent/.style={draw,
      -latex}]

      \node {Gerencia General \\ Biogás Colombia} child { node {Dirección de \\
            Operaciones (O\&M)} child { node {Supervisión SCADA y Control de Campo (GCCS)}
          } child { node {Mantenimiento de Motogeneradores y Skids} } } child { node
          {Dirección Administrativa \\ y Mercados} child { node {Gestión Regulatoria
        (UAESP / CREG / Vanti)} } child { node {Auditoría y Venta de Créditos (CER)} } }
        child { node {Dirección HSE y \\ Sostenibilidad} child { node {Laboratorio
        Químico y Control Emisiones} } };

  \end{tikzpicture}}
  \caption{Estructura organizacional detallada y jerarquía operativa de Biogás
  Doña Juana S.A.S.  E.S.P.}
  \label{fig:organigrama}
\end{figure}

\newpage

\section{Línea Base Estratégica}

La evaluación de la línea base estratégica requiere el análisis de la misión
corporativa, la visión operativa y su alineación con mandatos de política
pública y marcos de sostenibilidad.

\subsection{Misión Corporativa}

La misión de Biogás Doña Juana S.A.S.  E.S.P. se centra en la sostenibilidad
ambiental y la mitigación del cambio climático. La organización transforma el
biogás en un activo monetizable que mejora la calidad del aire local y
suministra energía renovable. Esta misión se alinea con los Objetivos de
Desarrollo Sostenible de las Naciones Unidas, específicamente los ODS 7 y 8, y
se operacionaliza en dos unidades de negocio principales.

La unidad de mitigación ambiental captura y destruye gases de efecto invernadero
para generar Reducciones Certificadas de Emisiones (CER) bajo el Mecanismo de
Desarrollo Limpio. Utilizando la metodología ACM0001, la organización registra
su producción en plataformas digitales y transacciona en mercados globales,
permitiendo a corporaciones compensar su huella de carbono.

La concesión exige un mecanismo de distribución de ingresos donde el 24\% de los
ingresos brutos por la venta de créditos de carbono y el 6\% de las ventas de
electricidad se transfieren a la UAESP. Estos fondos financian inversiones
sociales en las comunidades de Mochuelo y Usme, desarrollando infraestructura
pública y redes de saneamiento. La eficiencia en la generación de CER y las
transferencias financieras se documentan en la \cref{tab:cer_audit}.

\begin{table}[htbp]
  \centering
  \caption{Certificación histórica de CER y rendimientos financieros}
    \label{tab:cer_audit}
  \begin{tabular}{llrl}
    \toprule Periodo & Autoridad & CER Solicitados & Rendimiento a UAESP (COP)
    \\
    \midrule
    Sep 2009 - Dic 2009 & DNV & 76,048 & Consolidado anual \\
    Dic 2009 - May 2009 & DNV & 226,152 & Consolidado anual \\
    May 2010 - Sep 2010 & ICONTEC & 201,470 & Consolidado anual \\
    \midrule
    Total (Muestra) & Múltiple & 503,670 & $\sim$2.38 Mil Millones \\
    \bottomrule
  \end{tabular}
\end{table}

\subsection{Visión Corporativa}

La visión operativa exige maximizar la utilidad energética del biogás mediante
su integración en la matriz energética colombiana. El objetivo estratégico a
largo plazo es expandir la capacidad instalada a 15 megavatios (MW) para junio
de 2026. Esta expansión requiere la sincronización de nueve motogeneradores, lo
que permitirá a la entidad asegurar la solvencia financiera a través de la
comercialización de energía y reducir su dependencia de la volatilidad de los
mercados de carbono.

\newpage

\section{Formulación del Problema: La Paradoja de Capacidad, Interconexión y
Utilización}

El problema operativo fundamental de Biogás Doña Juana
S.A.S. E.S.P. es la asimetría entre la capacidad de captura de biogás
crudo y la generación efectiva de energía exportable a la red. Esta
discrepancia constituye un cuello de botella crítico, resultando en una
pérdida de ingresos corporativos y un rendimiento termodinámico subóptimo
para la matriz energética de Bogotá.

\subsection{Déficit de Recuperación Energética y Desperdicio Termodinámico}

La planta captura entre 13,500 y 16,000 metros cúbicos de biogás por hora. Los
análisis químicos establecen una concentración volumétrica media de 52.5\% de
metano (CH$_4$) y 38.1\% de dióxido de carbono (CO$_2$), junto con contaminantes
traza que incluyen 42 ppmV de sulfuro de hidrógeno (H$_2$S) corrosivo, siloxanos
y humedad. Esta concentración de metano constituye un vasto depósito continuo de
energía termodinámica.

Aunque la infraestructura de generación teórica proyectaba una capacidad
de hasta 25 MW, la realidad operativa diverge sustancialmente.
Actualmente, la planta suministra únicamente 1.7 MW de energía eléctrica
activa al sistema de interconexión nacional, un volumen suficiente para
abastecer a aproximadamente 4,000 hogares.

Esta subutilización crónica impone un paradigma operativo en el cual hasta
el 90\% del potencial energético del metano capturado debe ser quemado en
antorchas, en lugar de convertirse en electricidad o biometano purificado.
Si bien la combustión mitiga el impacto atmosférico del metano y permite
la generación de Reducciones Certificadas de Emisiones (CER), es un
proceso térmicamente destructivo que disipa la energía calórica intrínseca
del gas. Esto cumple el mandato de mitigación ambiental, pero fracasa en
la capitalización energética.

\subsection{Limitaciones Técnicas, Burocráticas e Infraestructurales}

Las causas fundamentales de esta paradoja de capacidad son multidimensionales,
abarcando limitaciones de ingeniería, déficits infraestructurales en la red
eléctrica externa y fricción burocrática con las entidades de control estatal.

\subsubsection{Cuello de Botella en la Interconexión Eléctrica}

Existe una restricción de infraestructura externa respecto a la capacidad de la
red de distribución municipal para absorber y transmitir la energía generada.
La integración de sistemas de generación distribuida en la red urbana principal
requiere subestaciones y líneas de transmisión de alto voltaje robustas.

Históricamente, la ausencia de infraestructura de red de alta capacidad en
las inmediaciones de Doña Juana, sumada a la falta de apoyo de las
autoridades energéticas municipales para modernizar las líneas de
transmisión, ha impedido que la planta exporte más de 1.7 MW. Generar y
despachar una carga continua de 15 MW requiere equipos de sincronización y
líneas especializadas que han permanecido incompletas o retrasadas por
fallas en la coordinación interinstitucional.

\subsubsection{Impurezas Químicas e Ingeniería de Alto Mantenimiento}

Desde una perspectiva mecánica, el biogás crudo derivado de residuos sólidos
urbanos contiene impurezas altamente abrasivas. El uso de una mezcla de gas con
42 ppmV de sulfuro de hidrógeno y siloxanos en motores de combustión interna de
trabajo continuo requiere instalaciones de pretratamiento y depuración química
de alta complejidad y mantenimiento.  Si los siloxanos no se eliminan,
precipitan como depósitos de sílice en las culatas, bujías y válvulas,
provocando fallas catastróficas en los motores. Escalar esta infraestructura de
purificación de 1.7 MW a 15 MW, utilizando nueve motogeneradores, exige un alto
gasto de capital. Esto requiere acuerdos de compra a largo plazo garantizados y
una interconexión de red estable para justificar la inversión.

\subsubsection{Fricción Burocrática y Vulnerabilidades Legales}

La relación entre el concesionario y la UAESP introduce una fricción legal que
retrasa las actualizaciones técnicas. La operación del relleno está sujeta a un
intenso escrutinio debido a su historial de emergencias ambientales y la presión
política para adoptar modelos de Basura Cero. En consecuencia, las operaciones
enfrentan auditorías de la Contraloría de Bogotá sobre la ejecución de contratos
gestionados por la UAESP. Además, el marco contractual ha generado disputas
legales ante el Consejo de Estado respecto a los límites de capacidad
volumétrica del relleno, específicamente la interpretación de la Resolución CAR
1351 de 2014, la cual delinea la capacidad base de 8.4 millones de metros
cúbicos y la Fase II de Optimización de 10.1 millones.

Esta incertidumbre legal y operativa disuade el gasto de capital corporativo. La
historia de la empresa subraya esta vulnerabilidad. En 2013, con pasivos de
aproximadamente 35,176 millones de COP frente a activos de 31,201 millones de
COP, Biogás Doña Juana S.A. fue sometida a un acuerdo de reorganización
empresarial supervisado por la Superintendencia de Sociedades para evitar su
liquidación. El acuerdo otorgó un período de gracia de 30 meses y un cronograma
de reestructuración de siete años, estabilizando la operación. Sin embargo, este
historial de patrimonio negativo demuestra que la empresa no puede sobrevivir a
largo plazo dependiendo exclusivamente de los ingresos volátiles del mercado de
carbono; es imperativo finalizar su capacidad de generación de energía para
asegurar una solvencia financiera permanente.

\newpage

\section{Demanda y Justificación: La Estrategia Energética de Doble Vector}

La justificación estratégica y económica para resolver la restricción de
capacidad y ejecutar la actualización de equipos a 15 MW, junto con la
exploración de vectores alternativos de utilización molecular, se fundamenta en
la demanda del mercado interno, imperativos ambientales urbanos e incentivos
fiscales macroeconómicos en Colombia.

\subsection{Vector de Demanda 1: Electricidad Renovable Firme}

La matriz energética colombiana, caracterizada por una baja intensidad de
carbono debido a su dependencia de la generación hidroeléctrica, es vulnerable a
fenómenos climáticos. El ciclo El Niño-Oscilación del Sur induce sequías que
reducen los niveles de los embalses y amenazan con el racionamiento.  En
consecuencia, existe una demanda nacional por Fuentes No Convencionales de
Energía (FNCE) firmes que garanticen energía de carga base despachable,
independiente de variaciones hidrológicas.

A diferencia de los recursos solares o eólicos intermitentes, la generación
eléctrica derivada del biogás opera de forma continua, constituyendo un activo
crítico para la estabilización de la red.  Suministrar 15 MW de potencia
continua a la red metropolitana de Bogotá mejora la resiliencia energética
urbana, reduce las pérdidas de transmisión y compensa la dependencia de la
generación térmica de respaldo. La modelación financiera demuestra que, mediante
incentivos a la inversión, los sistemas de biogás alcanzan Tasas Internas de
Retorno equivalentes a la generación diésel, pero generan un Valor Presente Neto
superior.

\subsection{Vector de Demanda 2: El Pivote de Interconexión de Biometano}

La demanda del mercado ha evolucionado hacia la utilización directa de la
molécula de gas purificada. Existe un mercado emergente para el Gas Natural
Renovable, o biometano. Transformar el biogás crudo en biometano de alta pureza
ofrece un rendimiento energético total que es teóricamente superior en un factor
de 2 a 3 respecto a su conversión en electricidad, al eludir las pérdidas
termodinámicas inherentes a los motores de combustión interna.

Biogás Colombia desarrolla un marco de asociación en este dominio. Vanti S.A.
ESP, distribuidor regional de gas natural, ha iniciado la planificación del
Proyecto RM-1-04718-2025 para diseñar y construir una red de interconexión
destinada a recibir biometano de la instalación Doña Juana.

\begin{table}[htbp]
  \centering
  \caption{Especificaciones técnicas para la inyección de biometano}
  \label{tab:biometano}
  \begin{tabular}{p{0.35\textwidth} p{0.55\textwidth}}
    \toprule
    Parámetro & Especificación Técnica \\
    \midrule
    Designación del proyecto & RM-1-04718-2025 \\
    Coordenadas & Longitud -74.147562, Latitud 4.492388 (1820m) \\
    Flujo proyectado & 2,580 $m^3$/h base, escalable a 3,400 $m^3$/h \\
    Presión operativa & Media 60 Psi. Alta 375-250 Psi \\
    Infraestructura & Tubería de polietileno de alta densidad (1 a 8 pulg.) \\
    Base de consumidores & 2,000 clientes en el Distrito 1 \\
    \bottomrule
  \end{tabular}
\end{table}

Inyectar biometano en la red urbana de Vanti apoya la economía circular al
descarbonizar el suministro municipal de gas. Permite a 2,000 clientes utilizar
moléculas renovables. Si la interconexión eléctrica de Enel Codensa persiste
como un cuello de botella, este gasoducto proporciona un vector de exportación
secundario de alta capacidad, permitiendo monetizar el gas sin requerir
actualizaciones en las líneas de transmisión eléctrica.

\subsection{Justificación Regulatoria, Fiscal y Financiera}

La viabilidad financiera para ejecutar el gasto de capital requerido para la
actualización eléctrica de 15 MW y la instalación de biometano está respaldada
por la transición del Estado colombiano hacia infraestructuras de energía
sostenible, codificada en la Ley 1715 de 2014 y ampliada por la Ley 2099 de
2021.

Estos instrumentos legislativos establecen un entorno fiscal para gastos de
capital en fuentes no convencionales de energía renovable, clasificando la
biomasa y el biogás como beneficiarios principales. Mediante un procedimiento
supervisado por la Unidad de Planeación Minero Energética (UPME) y la Autoridad
Nacional de Licencias Ambientales (ANLA), la empresa puede acceder a incentivos
fiscales que reducen el Costo Nivelado de la Energía (LCOE).

\begin{table}[htbp]
  \centering
  \caption{Marco de incentivos fiscales para gastos de capital en FNCE}
  \label{tab:incentivos}
  \begin{tabular}{llp{0.5\textwidth}}
    \toprule
    Mecanismo Fiscal & Base Normativa & Impacto Financiero \\
    \midrule
    Deducción de Renta & Ley 1715 (Art. 11) & Deducción de hasta el 50\% del
    valor de la inversión sobre la renta líquida. \\
    Depreciación Acelerada & Ley 1715 (Art. 14) & Tasa anual máxima del 33.33\%
    para maquinaria, equipos y obras civiles. \\
    Exclusión de IVA & Ley 1715 (Art. 12) & Exclusión total del IVA (19\%) en la
    compra nacional o importación de equipos vinculados. \\
    Exención Arancelaria & Ley 1715 (Art. 13) & Exención del pago de derechos
    arancelarios para la importación de maquinaria especializada. \\
    \bottomrule
  \end{tabular}
\end{table}

La modelación financiera demuestra que la aplicación acumulativa de estos
beneficios puede subsidiar el 46.1\% del gasto de capital inicial mediante
ahorros fiscales directos. Un modelo de inversión estandarizado para un gasto de
4,000 millones de COP proyecta 700 millones en ahorros de impuestos sobre la
renta, 760 millones en exención de IVA y un alivio del flujo de caja mediante
depreciación acelerada que supera los 385 millones.

Al complementar este soporte de capital con el flujo de ingresos generado por la
venta de reducciones certificadas (CER) en el mercado voluntario, frecuentemente
transados a tasas cercanas a 4.55 euros por tCO2e, la arquitectura financiera
para la expansión de 15 MW y la interconexión resulta robusta. Completar esta
expansión de doble vector alinea la operación con los objetivos jurídicamente
vinculantes del Plan de Acción Climática de Bogotá 2020-2050, el cual exige la
neutralidad de carbono y la integración de tecnologías de conversión de
residuos.

\newpage

\section{Fundamentación de Manufactura Flexible y Dinámica Molecular}

La concepción de la planta trasciende el simple tratamiento de residuos para
asimilarse a un sistema de manufactura flexible de fluidos. En esta
arquitectura, el flujo de biogás no es un desecho geológico, sino una variable
continua de energía química que debe ser enrutada dinámicamente según las
restricciones termodinámicas y topológicas del sistema. La flexibilidad de
este paradigma exige una comprensión profunda de las leyes físicas que
gobiernan la metrología del gas y una formulación matemática rigurosa para
su control de manera autónoma.

La medición precisa de las fracciones volumétricas de la mezcla fluida
requiere métodos ópticos y magnéticos que eludan la degradación progresiva de
los biosensores convencionales. La concentración de metano se determina
mediante espectrometría infrarroja no dispersiva, un fenómeno intrínsecamente
fundamentado en la atenuación óptica descrita por la ley de Beer-Lambert. La
absorción de la intensidad lumínica a través del medio obedece estrictamente
a la relación $I = I_{0} \exp(-\epsilon c L)$, donde $I_{0}$ es la intensidad
de referencia en el vacío, $\epsilon$ es la absortividad molecular del
enlace carbono-hidrógeno a $3.3~\mu\text{m}$, $c$ es la concentración local
y $L$ es la longitud del camino de interacción óptica. Para asegurar una
medida invariante ante deposiciones físicas en la cavidad geométrica, se
implementa una topología de canal de referencia óptico que anula
algebraicamente las perturbaciones exógenas del entorno. Concurrentemente,
la cuantificación del oxígeno disuelto aprovecha su susceptibilidad magnética
intrínseca. Al presentar un comportamiento paramagnético puro, la molécula
interactúa fuertemente con los gradientes de campos magnéticos aplicados.
Este fenómeno físico permite evaluar el flujo másico sin la necesidad
imperativa de utilizar elementos electroquímicos reactivos, los cuales
sufrirían una inactivación irreversible ante la presencia endémica de sulfuro
de hidrógeno.

El control topológico de la extensa red de extracción demanda un enfoque
predictivo de carácter decididamente multivariable. Para eludir las
representaciones matriciales convencionales y preservar la generalidad del
sistema, la dinámica del problema se formula naturalmente en el marco del
álgebra abstracta lineal. Sean un espacio de estados abstracto $X$ y un
espacio de acciones $U$, estructurados como espacios vectoriales normados
sobre el campo de los números reales. La dinámica ciberfísica busca, en cada
instante, minimizar un funcional de costo $J$, el cual se proyecta sobre un
horizonte temporal discreto. Este funcional adquiere la topología de una
suma de normas sobre el espacio producto
\begin{equation}
 J = \sum_{k} \|y(t+k|t) - r(t+k)\|_{Q}^{2} + \sum_{j} \|\Delta
 u(t+j|t)\|_{R}^{2},
\end{equation}

donde los mapeos $Q : X \to X$ y $R : U \to U$ constituyen operadores
lineales autoadjuntos y definidos positivos. Dichos operadores establecen
la estructura métrica del espacio funcional, penalizando implacablemente
cualquier divergencia del estado proyectado $y$ respecto a la trayectoria
ideal $r$, al tiempo que acotan la magnitud intrínseca de la acción de
control diferencial $\Delta u$. Esta arquitectura algorítmica evalúa las
múltiples trayectorias permitidas y mapea una respuesta que respeta
invariablemente las fronteras físicas de operación, aislando mecánicamente
cualquier discontinuidad o perturbación que induzca estrés anaeróbico en el
volumen de extracción.

La conversión termodinámica final de la corriente purificada exige procesos
de alta simetría para preservar el flujo de exergía. La exportación eléctrica
firme se ejecuta a través de generadores alternativos operando bajo el
principio de combustión pobre. En este régimen térmico, un margen masivo de
fluido inerte en la compresión reduce de forma exponencial la temperatura de
estancamiento en la cámara, imposibilitando la síntesis química indeseada de
óxidos de nitrógeno y garantizando una operación continua bajo los márgenes
de radiación permitidos. En paralelo, la reconfiguración para generar
biometano comercial se confía a la filtración mediante fibras capilares de
poliimida. Este mecanismo obedece teóricamente a un modelo riguroso de
solución y difusión macromolecular. Las especies polares como el dióxido de
carbono presentan una elevada constante de solubilidad en la matriz
dieléctrica asimétrica, atravesando la barrera polimérica de manera
inmediata. Como contrapartida, el metano se mantiene estructuralmente
aislado debido a su simetría tetraédrica perfecta, cuya neutralidad
electromagnética induce una profunda barrera cinética que lo confina a la
línea de suministro primario. Este elegante fraccionamiento permite derivar
la potencia química inherente del gas eludiendo las ineficiencias entrópicas
que rigen a las máquinas cíclicas de calor.

\newpage

\section{Ingeniería de Procesos y Ciclo de Producción}

El ciclo productivo de Biogás Doña Juana integra ingeniería de fluidos y
termodinámica para transformar un pasivo ambiental inestable en vectores
energéticos estandarizados. El proceso se articula en tres fases críticas:
extracción controlada, acondicionamiento químico y conversión energética.

\subsection{Fase 1: Sistema de Recolección y Control (GCCS)}

El Sistema de Recolección y Control de Gas (GCCS) actúa como el componente
respiratorio del vertedero. El proceso de producción inicia con la aplicación de
un campo de vacío mediante sopladores industriales que extraen activamente el
gas de los espacios porosos de la masa de residuos. Esta etapa requiere un
balance barométrico preciso gestionado por telemetría en tiempo real: un vacío
excesivo introduce oxígeno, eliminando las bacterias metanogénicas y creando
riesgo de incendio, mientras que un vacío insuficiente permite la fuga de
emisiones fugitivas a la atmósfera.

\subsection{Fase 2: Acondicionamiento y Eliminación de Siloxanos}

Antes de cualquier valorización, el gas crudo debe someterse a un pretratamiento
riguroso. El biogás ingresa a un sistema de acondicionamiento (skid) diseñado
para deshumidificar la mezcla y eliminar partículas. El paso crítico en esta
fase es la remoción de siloxanos, compuestos derivados de productos cosméticos
en descomposición. Sin este proceso, los siloxanos se convierten en sílice
(arena) dentro de las cámaras de combustión, causando daños abrasivos
irreversibles en pistones y válvulas. La eficiencia de este fregado químico es
determinante para garantizar la operatividad $24/7$ de la planta.

\subsection{Fase 3: Tecnologías de Valorización}

El gas acondicionado se dirige hacia dos líneas de producción distintas según el
producto final deseado:

\begin{itemize}
  \item \textbf{Generación Eléctrica
  (Motores Reciprocantes):} Para la producción de electricidad, el gas se
  inyecta en una flota de 9 motogeneradores, presumiblemente de la serie
  Jenbacher Tipo 4 o 6. Estos motores están endurecidos metalúrgicamente para
  soportar el poder calorífico variable del biogás y convertir la energía
  química en potencia eléctrica para la red local.
  \item \textbf{Producción de Biometano (Separación por Membranas):} Para la
    nueva línea de Gas Natural Renovable, la empresa implementa tecnología de
    separación por membranas. El biogás se comprime a alta presión y se fuerza a
    través de membranas selectivas donde las moléculas de $\text{CO}_2$ permean
    más rápido que las de $\text{CH}_4$. Este proceso físico, preferido por no
    requerir reactivos químicos, resulta en una corriente de biometano con
    pureza superior al 96\%.  Finalmente, el gas es comprimido entre 250 y 375
    psi para su inyección en el anillo de alta presión de Vanti.
\end{itemize}

\subsection{Oportunidades de Automatización y Control de Procesos}

El rendimiento termodinámico de la planta y la viabilidad de la transición hacia
la inyección de biometano con pureza superior al 96\% dependen intrínsecamente
de la calidad y estabilidad del biogás extraído. En este contexto, se identifica
una vulnerabilidad operativa en la dinámica de fluidos del campo, susceptible de
mejora mediante ingeniería de control avanzado.

\subsubsection{Problemática Identificada: Desbalance Barométrico y Dilución del
Biogás}

La dinámica de fluidos dentro de la matriz porosa de los residuos es altamente
heterogénea. Las variaciones en la presión barométrica superficial y la
subsidencia natural del terreno provocan fluctuaciones en el gradiente de
presión generado por los sopladores. Un vacío excesivo en zonas de baja densidad
introduce oxígeno atmosférico ($\text{O}_{2}$) a través de fisuras
superficiales. Físicamente, esto desencadena dos problemas críticos: la
inhibición competitiva de las bacterias metanogénicas (estrictamente anaerobias)
y la dilución de la concentración volumétrica de $\text{CH}_{4}$, reduciendo el
índice de Wobbe del gas.  Además, un nivel de $\text{O}_{2}$ superior al 3\%
incrementa el riesgo de ignición subterránea espontánea.

\subsubsection{Arquitectura del Sistema de Control e Instrumentación Actual}

Actualmente, la regulación se realiza mediante un control de supervisión
macro-estructural en la planta central, midiendo el flujo totalizado. A nivel de
instrumentación, la compañía cuenta con los siguientes sensores en su operación
base:
\begin{itemize}
  \item \textbf{Transmisores de presión y caudalímetros volumétricos:}
    Instalados exclusivamente en el cabezal principal de succión (antes de los
    sopladores) para monitorear el flujo totalizado de
    $13,500~\text{m}^3/\text{h}$.
  \item \textbf{Termopares industriales (tipo K o J):} Ubicados en las tres
    antorchas cerradas para monitorear y garantizar que la destrucción térmica
    del metano ocurra a temperaturas superiores a los $500~^{\circ}\text{C}$.
  \item \textbf{Analizadores portátiles de gases (ej.  GEM5000):} Equipos
    basados fundamentalmente en celdas electroquímicas de desgaste rápido, los
    cuales son utilizados por técnicos de campo en rondas físicas para realizar
    muestreos periódicos en los 290 cabezales de pozo.
\end{itemize}

Sin embargo, a nivel de la red de 290 pozos verticales, la regulación de presión
en los cabezales depende de válvulas de compuerta de accionamiento manual. Este
paradigma de control puramente reactivo adolece de tiempos muertos (\textit{dead
times}) considerables en el lazo de retroalimentación, imposibilitando una
respuesta dinámica ante caídas de presión atmosférica o cambios rápidos en la
permeabilidad del relleno.

\subsubsection{Mejora Planteada: Lazo de Control Predictivo basado en Sensores
IoT y NDIR}

Para estabilizar la fracción molar de metano y asegurar un flujo laminar y
constante hacia los 9 motogeneradores proyectados, se propone la implementación
de un sistema de Control Predictivo Basado en Modelos (MPC) o un control PID en
cascada, instrumentado a nivel de \textit{manifolds} (sub-colectores) que
agrupan conjuntos de 10 a 15 pozos.

El plan de automatización se estructura en las siguientes capas tecnológicas:

\begin{itemize}
  \item \textbf{Capa de Instrumentación Analítica (Sensores):} Instalación de
    instrumentación \textit{in-situ} en cada manifold, seleccionada
    estratégicamente por su robustez:
    \begin{itemize}
      \item \textbf{Sensores Infrarrojos No Dispersivos (NDIR) para
        $\text{CH}_{4}$ y $\text{CO}_{2}$:} Operan bajo la ley de Beer-Lambert,
        midiendo la atenuación de la radiación infrarroja ($3.3~\mu\text{m}$
        para el metano). \textbf{Justificación:} A diferencia de las celdas
        electroquímicas, permiten mediciones continuas sin desgaste de
        reactivos, garantizando alta vida útil operativa.
      \item \textbf{Celdas Paramagnéticas para $\text{O}_{2}$:} Aprovechan la
        susceptibilidad magnética positiva del oxígeno.  \textbf{Justificación:}
        Las celdas electroquímicas tradicionales se envenenan y fallan al
        exponerse al $\text{H}_{2}\text{S}$ del biogás.  Las paramagnéticas son
        físicamente inmunes a esta interferencia, previniendo riesgos de
        explosión.
      \item \textbf{Transmisores de Presión Diferencial Piezorresistivos:}
        \textbf{Justificación:} Proveen la variable de proceso inmediata para
        que el algoritmo detecte fluctuaciones de vacío local.
    \end{itemize}

  \item \textbf{Capa de Actuación Proporcional:} Reemplazo de las válvulas
    manuales en los \textit{manifolds} por válvulas de mariposa de control
    proporcional con caracterización isoporcentual, equipadas con actuadores
    electroneumáticos inteligentes. Estos recibirán señales estandarizadas de
    4-20 mA o comandos digitales.

  \item \textbf{Capa de Control Lógico y SCADA:} Los sensores transmitirán
    variables de proceso continuas mediante protocolos industriales (ej. Modbus
    TCP/IP o Profinet) hacia Controladores Lógicos Programables (PLC) en
    topología de anillo de fibra óptica.  El algoritmo de control priorizará la
    calidad del gas:
    \begin{equation}
      u(t)=K_{p}e(t)+K_{i}\int_{0}^{t}e(\tau)d\tau+K_{d}\frac{de(t)}{dt}.
    \end{equation}
    Donde la variable de proceso primaria será la concentración de $\text{CH}_{4}$ y
    la restricción estricta (override control) será mantener el
    $\text{O}_{2}<1,5\%$.
\end{itemize}

Si se detecta un pico de oxígeno por intrusión atmosférica, el lazo de control
estrangulará automáticamente la válvula del \textit{manifold} correspondiente
para disminuir el vacío, aislando la perturbación sin afectar la red global.
Esta automatización permite a la red de polietileno de 76 km comportarse como un
sistema termodinámico autorregulado, minimizando la intervención humana en las
623 hectáreas y garantizando la calidad del combustible requerida para el
despacho firme de los 15 MW.

\newpage

\section{Evaluación Estratégica: Análisis FODA Integral}

Un diagnóstico organizacional exhaustivo requiere un análisis FODA.  Esta matriz
evalúa las capacidades operativas internas de Biogás Doña Juana frente a las
condiciones del mercado y la regulación externa, definiendo su postura
estratégica de cara a la expansión de capacidad proyectada para junio de 2026.

\begin{table}[htbp]
  \centering
  \caption{Matriz analítica FODA de Biogás Doña Juana S.A.S. E.S.P.}
  \label{tab:foda}
  \begin{tabular}{p{0.45\textwidth} p{0.45\textwidth}}
    \toprule
    Fortalezas & Oportunidades \\
    \midrule
    Red de extracción de biogás extensa (290 pozos, 76 km de tubería HDPE).
    \newline
    Experiencia técnica en tratamiento anaeróbico y quema (13,500 $m^3$/h).
    \newline
    Dominio en mercados de compensación de carbono.
    \newline
    Concesión a 35 años que asegura el monopolio de la materia prima.  &
    Expansión de la red a 15 MW con 9 motogeneradores para junio de 2026.
    \newline
    Proyecto de inyección de biometano (RM-1-04718-2025) con Vanti.
    \newline
    Incentivos fiscales (Ley 1715/2099) que subsidian la
    inversión de capital. \newline Alineación con los
    planes de acción climática 2050 de Bogotá. \\
    \midrule
    Debilidades & Amenazas \\
    \midrule
    Incapacidad de transformar el flujo volumétrico en energía (1.7 MW frente a
    15 MW).
    \newline
    Dependencia operativa de la infraestructura de red externa (Enel Codensa).
    \newline
    Historial de volatilidad financiera (reorganización en 2013).
    \newline
    Alto costo de mantenimiento por impurezas del gas (siloxanos, H2S). &
    Fricción legal con la UAESP sobre los límites de la concesión (Resolución
    CAR 1351).
    \newline Presión sociopolítica para erradicar el modelo de relleno sanitario.
    \newline Auditorías de la Contraloría sobre el modelo de ingresos compartidos.
    \newline Volatilidad de precios en el mercado de créditos de carbono (CER).
    \\
    \bottomrule
  \end{tabular}
\end{table}

\subsection{Análisis de Capacidades Internas}

La empresa posee una infraestructura de extracción robusta. La gestión continua
de 76 kilómetros de tubería interna y 290 pozos dinámicos demuestra el control
técnico sobre la topografía inestable del relleno. Esta red asegura un flujo de
13,500 $m^3$/h, garantizando el suministro para la generación eléctrica o la
purificación de biometano. Asimismo, la validación bajo estándares
internacionales (CDM ACM0001) permite la comercialización de Reducciones
Certificadas de Emisiones (CER), generando un activo financiero líquido mediante
asociaciones corporativas.

La debilidad principal es la ineficiencia en la conversión energética. La
necesidad de quemar térmicamente hasta el 90\% del metano capturado por falta de
capacidad de generación sincronizada representa una pérdida de valor de mercado.
Suministrar únicamente 1.7 MW frente a un objetivo de 15 MW implica un costo de
oportunidad sustancial. Adicionalmente, el historial corporativo revela
vulnerabilidades financieras; en 2013, la empresa enfrentó una reorganización
debido a un patrimonio negativo de 3.966 mil millones de COP frente a pasivos de
35.176 mil millones de COP. Esto subraya el riesgo estructural de depender
exclusivamente de los mercados de carbono sin ingresos estabilizados por venta
de energía. Finalmente, la operación depende de la modernización de las
subestaciones de Enel Codensa para permitir el despacho a la red nacional.

\subsection{Análisis del Entorno Externo}

El entorno macroeconómico y regulatorio ofrece dos vías de expansión: la
generación eléctrica a escala y la distribución de biometano. Alcanzar el
objetivo de 15 MW para 2026 transformará la estructura de ingresos de la
empresa, convirtiéndola en un proveedor de servicios públicos de carga base.

Simultáneamente, el proyecto de biometano con Vanti permite purificar el gas
para su inyección en la red regional a 60 Psi, con una capacidad de hasta 3,400
$m^3$/h. Esta alternativa de exportación molecular mitiga el riesgo de
restricciones en la transmisión eléctrica. La ejecución de estos proyectos es
económicamente viable gracias a los escudos fiscales provistos por la UPME bajo
la Ley 1715.

La viabilidad a largo plazo está condicionada por la continuidad operativa del
Relleno Sanitario Doña Juana. El modelo de disposición final enfrenta presión
política orientada hacia paradigmas de cero residuos. Una alteración en la
política municipal que reduzca el flujo de residuos orgánicos agotaría
inevitablemente la materia prima de la planta. Adicionalmente, el marco de la
concesión genera fricción legal con la UAESP respecto a los límites volumétricos
definidos por las licencias ambientales, específicamente la Resolución CAR 1351
de 2014. La administración de las transferencias contractuales (24\% de ingresos
por CER y 6\% por electricidad) está sujeta a auditorías rigurosas de la
Contraloría de Bogotá, lo cual introduce fricción operativa recurrente.

\newpage

\section{Síntesis y Perspectiva Estratégica}

Biogás Doña Juana S.A.S. E.S.P. ocupa una posición contradictoria en la
infraestructura ambiental de Bogotá. La organización ejecuta la mitigación de
gases de efecto invernadero gestionando un entorno subterráneo de 623 hectáreas
para extraer y destruir contaminantes, lo que genera valor líquido en los
mercados internacionales de compensación de carbono. Sin embargo, como productor
de energía, la empresa presenta un subdesarrollo crítico. La discrepancia entre
el volumen de energía termodinámica capturada y la energía comercializada
constituye su desafío estratégico central.

El imperativo corporativo de financiar y ejecutar la actualización de la
capacidad de generación a 15 MW, utilizando nueve motogeneradores para junio de
2026, es indispensable para la viabilidad a largo plazo. La dependencia de la
quema térmica es insostenible desde la perspectiva del costo de oportunidad
económico y la optimización macroeconómica de recursos.

El entorno legislativo colombiano proporciona condiciones óptimas para esta
capitalización. Los beneficios fiscales de las Leyes 1715 de 2014 y 2099 de
2021, que incluyen deducciones de renta del 50\%, depreciación acelerada y
exenciones de IVA y aranceles, establecen una vía estructurada para la
adquisición de equipos de generación y purificación.

Para superar los cuellos de botella técnicos y burocráticos, la organización
debe implementar una estrategia de producción energética de doble vector.
Mientras negocia con las empresas de servicios públicos para asegurar las
actualizaciones de transmisión de alto voltaje necesarias para absorber la carga
de 15 MW, la entidad debe acelerar el proyecto de integración del gasoducto de
biometano con Vanti (RM-1-04718-2025). La purificación del biogás crudo a gas
natural renovable para inyección directa en la red municipal elude la congestión
de la red eléctrica, estableciendo un flujo de ingresos secundario que atiende
hasta 2,000 clientes directos.

La transición de un proyecto de remediación ambiental a una empresa de servicios
energéticos diversificada requiere la gestión rigurosa de sus variables
operativas y regulatorias. La empresa debe mantener la contabilidad de carbono
para sostener los ingresos por CER, colaborar con la UAESP y la Contraloría para
asegurar la reinversión de dividendos sociales en las comunidades de Mochuelo y
Usme, y aplicar los incentivos de política pública a su estructura de capital.

Al resolver la paradoja de capacidad e interconexión mediante el suministro de
electricidad e inyección de biometano, Biogás Doña Juana completará la
transición de la mitigación de una crisis ambiental a la provisión energética
sostenible del Distrito Capital.

\newpage

\printbibliography

\nocite{*}

\end{document}
